\documentclass[11pt]{article}

\usepackage[margin=1in]{geometry}
\usepackage{booktabs}
\usepackage{caption}
\usepackage{siunitx}

\sisetup{
  round-mode = places,
  round-precision = 9,
  table-number-alignment = center
}

\title{SCIP Objective Values for the Fractional Partition MINLP}
\author{}
\date{}

\begin{document}

\maketitle

\begin{table}[htbp]
  \centering
  \caption{Objectives returned by the compact SCIP formulation}
  \label{tab:fractional-partition-objectives}
  \begin{tabular}{
    c
    S[table-format=1.9]
    S[table-format=1.9]
    S[table-format=1.9]
  }
    \toprule
    {$m$} & {$a=\frac{1}{2}$} & {$a=\frac{1}{4}$} & {$a=\frac{1}{100}$} \\
    \midrule
    3 & 1.333333 & 1.586934 & 5.586328 \\
    4 & 1.325694 & 1.578781 & 5.606681 \\
    5 & 1.329150 & 1.561089 & 5.58994 \\
    6 & 1.333333 & 1.586934 & 5.586328 \\
    \bottomrule
  \end{tabular}
\end{table}

\noindent
The runs used \(b=1\), \(0\leq l_i\leq 30\), group and denominator
tolerance \(10^{-5}\), and a limit of \(50000\) branch-and-bound nodes.
The auxiliary objective variable used the valid upper bound
\[
  r \leq 1+\frac{1}{6a}.
\]
The cases \((m,a)=(3,\frac12)\) and \((6,\frac12)\) ended with SCIP
status \texttt{optimal}. Every other run ended with status
\texttt{nodelimit}; those entries are the best feasible objectives found
within the node limit.

\end{document}
